%======================================================================
\section{Introduction}
\label{sec:introduction}
%======================================================================

Ensemble filters estimate the uncertainty of a forecast from a small
ensemble of model runs, and in high-dimensional systems they work only
because of localization: each region is analysed from nearby observations
with covariances estimated from the ensemble, and distant, spurious
correlations are suppressed \citep{houtekamermitchell2001,ott2004,hunt2007}.
The case for it is strongest in spatially extended chaos. There the forecast
errors within a region span few directions \citep{patil2001}, while the
number of unstable directions of the whole attractor, like its
Kaplan--Yorke dimension $\DKY$, grows in proportion to the domain
\citep{ruelle1982,manneville1985,egolfgreenside1994}. A filter without
localization needs at least as many members as there are unstable and
neutral directions \citep{bocquetcarrassi2017,takeda2026}, so its ensemble
must grow with the domain; a localized one needs only enough members for
the directions inside its local region.

Data assimilation in a learned latent space has so far given up this
remedy. Filters in the latent space of an autoencoder compress the whole
field into one latent vector and estimate its covariance from the ensemble
\citep{peyron2021,chen2023roadenkf,pasmans2025,tong2026laeenkf,si2025latentensf}.
A latent vector has no notion of distance, so there is nothing to taper:
recent latent filters either do not localize or leave localization for
future work \citep{chen2023roadenkf,tong2026laeenkf}, and the dimension of
the vector, fixed at training time, ties the model to one domain size. For
the large, extensively chaotic domains on which ensemble data
assimilation is used in practice, both are limits of principle rather than
of implementation.

A companion paper \citep[hereafter Paper~I]{jones2026transfer} makes the latent a
\emph{field}: a few channels at each site of a lattice that tiles the
domain, encoded, decoded and propagated by operators that are shared by
all sites and read a bounded neighbourhood. One trained model therefore
runs on a domain of any size by adding sites. Paper~I validates this
transfer against the model's own Lyapunov spectrum: trained on a small periodic
box or on windows of one large simulation, Paper~I's two-dimensional Kuramoto--Sivashinsky (KS)
models reproduce the true $\DKY$ to within $3.5\%$ at $44^2$ and $66^2$ (this paper's noisy-trained model:
$+3.6\%$ and $-2.1\%$, Table~\ref{tab:model_validation}), and
Kolmogorov flow needs the window route. A latent field has a geometry, and this paper asks
what that geometry is worth to an ensemble filter.

The answer has three parts. First, locality can be made exact: with a
decoder that reconstructs each patch from its own site, the observations
of a patch depend on that site alone, the latent observation operator is
block-diagonal, and a per-site ensemble filter in the latent space
needs no truncation of it. Second, the error that exact locality
leaves, tied to the patch lattice, is removed by running the filter on
shifted copies of the lattice and blending the analyses, without
retraining. Third, the localized filter keeps its accuracy and its
ensemble size as the domain grows, because both are set by the local
problem. These are properties of the method. Whether to use a learned
model at all is a separate question, because a filter built on the
correct equations is far more accurate than one built on a learned
surrogate. Models of real systems are not correct, however, and for many
systems, convection-permitting models of the atmosphere among them, the
errors are structural and poorly known, while the errors of the
observations are well characterized. We therefore train the model only on
noisy fields and compare it with physics-based filters whose models are
wrong in controlled ways and which are given the standard remedies for
model error.

\para{Contributions.}
\begin{itemize}
\item \emph{Exact locality in a learned latent space}
  (Sec.~\ref{sec:method}). With a per-patch decoder the latent observation
  Jacobian is block-diagonal (Proposition~\ref{prop:blockdiag}) and is
  obtained from as many forward passes as there are channels, at any
  domain size.
\item \emph{Shifted tilings} (Sec.~\ref{sec:da_tilings}). Four copies of the
  patch lattice shifted by half a patch, blended with smooth weights, lower
  the analysis error by about $30\%$ with the same model; the gain comes
  from the offsets, not from the extra members.
\item \emph{One model and one setting from $44^2$ to $440^2$}
  (Sec.~\ref{sec:da_2d}). In 2D KS the analysis error is $0.30$--$0.33$ on
  every domain over a factor of 100 in area, against $0.81$--$0.83$ for
  interpolating the observations, with an ensemble that saturates at
  about 32 members per tiling at both sizes tested, while an unlocalized filter
  fails once the attractor has more unstable directions than the ensemble
  has members.
\item \emph{What the nonlinear decoder costs} (Sec.~\ref{sec:da_nonlinear}).
  Observed through the decoder, the latent observation operator is
  nonlinear even for linear observations, and the filter linearizes it. At
  the native configuration the neglected term is about $4\%$ of the
  observation error, and relinearizing the update changes the error by at
  most $0.001$. With the true equations as forecast and the members in the
  decoder's range, the filter is no better than with the learned model; the
  error comes from the decoder's range, which costs about half of it, not
  from the learned dynamics or the nonlinearity.
\item \emph{Learned versus physics-based filters}
  (Sec.~\ref{sec:da_model_error}). The latent filter loses to a
  physics-based filter whose model error lies in a parameter or a bias,
  at every magnitude tested, once online parameter estimation is working, and wins when
  the error lies in the form of the equations and the part bias correction
  cannot remove exceeds about a quarter to a third of the observation error after the remedy,
  in the configuration tested. A missing process of
  increasing strength, tuned for on both sides, gives a tie just below that
  point and wins of $1.27$ and $1.29$ times above it.
\item \emph{Training on noisy data} (Sec.~\ref{sec:da_noisy}). The model is
  trained only on complete fields that carry observation-level noise, and
  assimilates at least as well as the same model trained on clean fields
  ($4$--$5\%$ lower error at every size, below the $10\%$ of our decision rule).
\end{itemize}

All experiments are twin experiments. The main system is 2D KS
(Sec.~\ref{sec:system}); 1D KS shows how the ensemble an unlocalized filter
needs grows with the domain and that a distance taper beats a distance-free
correction (Sec.~\ref{sec:da_1d}), and Kolmogorov flow shows that the
filter carries over to chaos organized by a forcing
(Sec.~\ref{sec:kolmogorov_da}). The target we have in mind is the
assimilation of convection-permitting simulations of organized convection,
where the equations of the cloud processes are uncertain and the domain is
set by the application. Section~\ref{sec:background} reviews the
background, Sec.~\ref{sec:method} the model and the filter, and
Secs.~\ref{sec:da_results}--\ref{sec:kolmogorov_da} the results;
Sec.~\ref{sec:discussion} discusses them and their limitations.
